\chapter{Addiction Recovery}

\section*{Definition and Overview}
Addiction recovery is the ongoing process by which a person with addiction achieves and maintains a life free from harmful substance use. It is much more than simply stopping the drug: recovery involves rebuilding health, relationships, and purpose, and it continues long after the last use. Recovery is different for everyone, and it may mean total abstinence or controlled, problem-free use, but it is generally described as a process of change rather than a single event.

Recovery is best understood as a personal and social journey supported by medical care, therapy, and community. It is a realistic and achievable outcome for most people, even after years of addiction or several earlier failed attempts.

\section*{Epidemiology}
Recovery is common. Large surveys suggest that many people who once met the criteria for addiction no longer do so, and most people who stop using do so without formal treatment. Relapse, however, is also common, and many people in recovery have made several earlier attempts.

Because addiction can be a lifelong condition, the number of people who are currently in recovery at any given time is difficult to measure precisely. What is clear is that recovery happens at all ages and across all social backgrounds, and that sustained recovery is more likely with ongoing support.

\section*{Aetiology and Pathophysiology}
Recovery begins with a change in the brain as much as a change in behaviour. With abstinence, the over-stimulated reward system gradually resets, and the circuits involved in decision-making and impulse control slowly strengthen. This process, sometimes described by the term neuroplasticity, takes weeks to months, and it is one reason that cravings and relapse risk can persist for a long time.

Recovery is supported when triggers, stress, and underlying mental health conditions are addressed at the same time. The brain changes of addiction do not simply reverse overnight; they require repetition, structure, and patience, much like building any new skill.

\section*{Clinical Features}
Recovery is not a set of symptoms, but certain features are characteristic:

\begin{itemize}
    \item \textbf{Abstinence or controlled use:} a period of reduced or absent substance use
    \item \textbf{Reduced cravings:} urges become less frequent and intense over time
    \item \textbf{Physical restoration:} rebuilding of sleep, appetite, and energy
    \item \textbf{Re-engagement:} renewed involvement in family, work, and hobbies
    \item \textbf{New coping skills:} developing strategies for managing stress and cravings
    \item \textbf{Lapses:} occasional slips, especially in the early stages of recovery
\end{itemize}

\section*{Differential Diagnosis}
Not all change looks the same, and it is useful to distinguish between several situations:

\begin{itemize}
    \item A lapse (a one-off use) from a full relapse into old patterns
    \item Recovery without formal treatment from recovery with professional support
    \item Abstinence due to illness, incarceration, or pregnancy from a genuine change of direction
    \item Ongoing mental illness from the emotional ups and downs of early recovery
\end{itemize}

\section*{When to Seek Medical Care}
Anyone in recovery who feels at risk of relapse, or who has had a relapse, should seek support promptly; the sooner help is sought, the easier it is to get back on track. Medical care is also needed for new or worsening mental health problems, or if medicines used to support recovery require review.

\textbf{Call 000 (or local emergency services) for:}
\begin{itemize}
    \item A relapse followed by overdose symptoms: drowsiness, slowed breathing, or collapse
    \item Suicidal thoughts or self-harm
    \item Severe withdrawal symptoms after a relapse and re-use
    \item Confusion, seizures, or loss of consciousness
\end{itemize}

\section*{Diagnosis and Investigations}
Recovery is monitored rather than formally diagnosed:

\begin{itemize}
    \item \textbf{Regular review:} check-ins with a GP, counsellor, or case worker
    \item \textbf{Substance testing:} urine or breath testing may be part of treatment programmes
    \item \textbf{Blood tests:} liver function and other checks to monitor damage from past use
    \item \textbf{Mental health review:} screening for depression, anxiety, and post-traumatic stress
    \item \textbf{Relapse prevention planning:} identifying triggers and rehearsing coping responses
\end{itemize}

\section*{Management}
Recovery is supported by a range of services that can be combined:

\begin{itemize}
    \item \textbf{Aftercare programmes:} ongoing counselling or structured follow-up after detoxification or rehabilitation
    \item \textbf{Mutual support groups:} such as Alcoholics Anonymous, Narcotics Anonymous, and SMART Recovery
    \item \textbf{Medicines:} such as naltrexone or buprenorphine to reduce craving and prevent relapse
    \item \textbf{Psychological therapy:} cognitive behaviour therapy to change thinking and behaviour
    \item \textbf{Family and social support:} involving partners, family, and friends in the process
    \item \textbf{Practical support:} help with housing, employment, and finances
    \item \textbf{Relapse prevention skills:} learning to handle triggers, stress, and high-risk situations
\end{itemize}

\section*{Complications}
\begin{itemize}
    \item Relapse, and the danger of overdose if tolerance has fallen
    \item Loss of motivation and the return of cravings in early recovery
    \item Unresolved or worsening mental health problems
    \item Relationship and financial strain during the recovery period
    \item Guilt, shame, or social isolation
    \item Health complications from years of past substance use
\end{itemize}

\section*{Prognosis and Outcomes}
Recovery is achievable. Many people who enter recovery maintain it for years, and some go on to live completely free of substance problems. Outcomes improve with longer engagement in treatment, good social support, and management of underlying mental health issues.

Relapse is not a sign of failure but a common part of the journey, and most people who relapse eventually achieve stable recovery. The strongest predictor of lasting success is sustained participation in treatment and support over months and years, rather than any single intervention.

\section*{Prevention and Self-Care}
\begin{itemize}
    \item Stay connected with supportive people and recovery groups
    \item Identify personal triggers and plan ahead for high-risk situations
    \item Build a daily routine with sleep, meals, exercise, and meaningful activity
    \item Learn and practise stress-management and relaxation skills
    \item Take medicines as prescribed and attend all appointments
    \item Treat mental health problems early rather than waiting
    \item Celebrate progress, and seek help quickly after a lapse
\end{itemize}

\section*{Sources and Further Reading}
The following organisations provide reliable information on addiction recovery and relapse prevention.

\begin{itemize}
    \item NHS. \textit{Information on addiction recovery and support services.} https://www.nhs.uk/conditions/
    \item World Health Organization. \textit{Resources on substance use and mental health.} https://www.who.int/
    \item NICE. \textit{Clinical guidance on drug and alcohol treatment and aftercare.} https://www.nice.org.uk/
    \item Mayo Clinic. \textit{Guide to addiction recovery and relapse prevention.} https://www.mayoclinic.org/
    \item MedlinePlus. \textit{Health information on substance use and recovery.} https://medlineplus.gov/
\end{itemize}
